\section{Learning-Augmented Influence Maximization}
\label{sec:method}

\subsection{Problem and Learned-Advice Interface}
Let $G=(V,E)$ denote a diffusion graph and let $\sigma(S)$ be the expected spread of seed set $S\subseteq V$. Given a seed budget $k$, influence maximization seeks
\begin{equation}
    \max_{S\subseteq V,\, |S|\le k} \sigma(S).
\end{equation}
At a sequential state $S$, the conditional marginal gain of candidate $v\notin S$ is
\begin{equation}
    \Delta(v\mid S)=\sigma(S\cup\{v\})-\sigma(S).
\end{equation}
We assume access to a learned advice model $\widehat{\Delta}(v\mid S)$ that is inexpensive relative to a high-fidelity classical influence oracle. Crucially, we make no assumption that this advice is uniformly accurate. The learned model proposes where expensive computation should be spent; an audited classical layer retains control of the actual seed decision. Figure~\ref{fig:overview} summarizes this interface.

\begin{figure}[t]
    \centering
    \resizebox{0.98\linewidth}{!}{\begin{tikzpicture}[
  font=\scriptsize,
  box/.style={draw, rounded corners=2pt, align=center, minimum height=9mm, inner sep=2.5pt},
  main/.style={box, text width=20mm},
  small/.style={box, text width=23mm},
  flow/.style={-{Latex[length=2mm]}, semithick},
  back/.style={-{Latex[length=2mm]}, semithick, dashed}
]
\node[main] (graph) {Full graph $G$\\$15{,}233$ nodes};
\node[main, right=4mm of graph] (screen) {Cheap global\\RR/RIS screening};
\node[main, right=4mm of screen] (short) {Compact shortlist\\$C_M$, $M=128$};
\node[main, right=4mm of short] (learn) {State-aware advice\\$\widehat{\Delta}(v\mid S)$};
\node[main, right=4mm of learn] (audit) {Predicted head +\\sentinel audit};

\node[small, below=9mm of audit] (fallback) {Expand oracle /\\classical fallback};
\node[small, left=5mm of fallback] (verify) {Progressive candidate +\\sample verification};
\node[small, left=9mm of verify] (select) {Select verified winner\\and update $S$};

\draw[flow] (graph) -- (screen);
\draw[flow] (screen) -- (short);
\draw[flow] (short) -- (learn);
\draw[flow] (learn) -- (audit);
\draw[flow] (audit) |- node[pos=0.30, right] {unsafe} (fallback);
\draw[flow] (audit) |- node[pos=0.30, left] {safe} (verify);
\draw[flow] (fallback) -- (select);
\draw[flow] (verify) -- (select);
\draw[back] (select.north) to[out=110,in=250] node[midway, left] {$|S|<k$} (learn.south);
\node[align=center, font=\scriptsize\bfseries, above=4mm of short] {Learned advice allocates expensive oracle work; audited verification retains decision control.};
\end{tikzpicture}
}
    \caption{Overview of the learning-augmented IM pipeline. Independent full-graph screening creates a high-recall shortlist; a state-aware predictor proposes promising candidates; residual auditing decides whether progressive verification is sufficient or classical fallback is required.}
    \label{fig:overview}
\end{figure}

\subsection{Full-Graph Candidate Screening}
The full graph is first reduced to a high-recall shortlist $C_M$ using a cheap global RR/RIS sketch. This separates global scalability from the harder sequential conditional-ranking problem. Screening is deliberately independent of the learned marginal model: its role is only to retain plausible strong seeds before the state-dependent stage. Unless otherwise specified, our default shortlist size is $M=128$, while $M\in\{64,128,256\}$ is used to expose the quality--cost sensitivity.

\subsection{State-Aware Marginal Proposal}
For candidates in $C_M\setminus S$, a graph predictor estimates $\widehat{\Delta}(v\mid S)$. Training explicitly exposes candidate nodes under multiple seed states and supervises conditional changes in gain. This matters because a predictor can otherwise learn a candidate-strength shortcut: it assigns nearly static scores to intrinsically strong nodes while missing overlap-induced diminishing returns. We therefore treat cross-state sensitivity of the same candidate as a first-class diagnostic rather than relying only on aggregate regression accuracy.

\subsection{Audited Residual Trust Test}
At each decision step, we evaluate a small predicted head together with rank-spaced sentinel candidates using a limited-budget classical oracle. For each audited candidate, define the residual
\begin{equation}
    r(v\mid S)=\Delta_{\mathrm{audit}}(v\mid S)-\widehat{\Delta}(v\mid S).
\end{equation}
The audited residuals estimate whether an unverified outsider can plausibly overturn the current audited winner. The trust statistic is therefore aligned with the downstream failure event that matters: excluding a candidate that could beat the verified head. When the audit is inconclusive, the method spends more oracle work rather than accepting the learned top candidate by default.

\subsection{Progressive Verification and Fallback}
When the audit indicates that the current head is safe, the solver progressively increases candidate coverage and influence-estimation budget only as needed. If the residual audit signals unreliable advice, verification expands; in the limiting case the method falls back to all-candidate classical evaluation. Consequently, informative advice can reduce expensive evaluations while sufficiently poor advice drives computation back toward the classical solver.

\begin{center}
\fbox{\begin{minipage}{0.94\linewidth}
\small
\textbf{Algorithm 1: Audited learning-augmented influence maximization.}\\[-1mm]
\textbf{Input:} graph $G$, budget $k$, shortlist size $M$, learned advice $\widehat{\Delta}$, classical influence oracle.\\[0.5mm]
\textbf{1. Screen.} Build a high-recall shortlist $C_M$ from the full graph using an independent RR/RIS sketch.\\
\textbf{2. Propose.} At state $S$, score $v\in C_M\setminus S$ with $\widehat{\Delta}(v\mid S)$ and form a predicted head plus rank-spaced sentinels.\\
\textbf{3. Audit.} Evaluate the audit set with a small classical budget and estimate residual uncertainty relative to the current audited winner.\\
\textbf{4. Verify progressively.} If outsider risk is controlled, increase candidate coverage and sample budget only as needed to verify the next decision.\\
\textbf{5. Fall back when needed.} If the audit remains unsafe, expand classical evaluation up to the full candidate set.\\
\textbf{6. Update.} Add the verified winner to $S$ and repeat Steps 2--5 until $|S|=k$.
\end{minipage}}
\end{center}
